Claim:
$\sqrt[3]{413 \mspace{1mu}} > 7 + \dfrac{4}{9}$.
This can be shown as follows:
\begin{align*}
  \paren{7 + \frac{4}{9}}^3
  &= \, 7^3 + 3 \cdot 7^2 \cdot \frac{4}{9} + 3 \cdot 7 \cdot \frac{4^2}{9^2} + \frac{4^3}{9^3} \,
   = \, 343 + \frac{196}{3} + \frac{112}{27} + \frac{64 / 27}{27} \, = \\
  &= \, 343 + 65 + \frac{1}{3} + 4 + \frac{4}{27} + \frac{2 + 10 / 27}{27} \,
   < \, 412 + \frac{1}{3} + \frac{9}{27} \, = \\
  &= \, 412 + \frac{2}{3}
   < 413
\end{align*}

Claim:
$6 \mspace{2mu} + \sqrt[3]{3 \mspace{1mu}} < 7 + \dfrac{4}{9}$;
that is:
$\sqrt[3]{3 \mspace{1mu}} < 1 + \dfrac{4}{9}$.
This can be shown as follows:
\begin{align*}
  \paren{1 + \frac{4}{9}}^3
  &= \, 1 + 3 \cdot \frac{4}{9} + 3 \cdot \frac{4^2}{9^2} + \frac{4^3}{9^3} \,
   = \, 1 + \frac{4}{3} + \frac{16}{27} + \frac{64}{729} \,
   = \, 2 + \frac{1}{3} + \frac{9 + 7}{27} + \frac{64 / 27}{27} = \\
  &= \, 2 + \frac{1}{3} + \frac{1}{3} + \frac{7}{27} + \frac{2 + 10 / 27}{27} \,
   = \, 2 + \frac{2}{3} + \frac{9}{27} + \frac{10}{729}
   > 3
\end{align*}

So, in conclusion:
\begin{equation*}
  6 \mspace{2mu} + \mspace{-1mu} \sqrt[3]{3 \mspace{1mu}} \mspace{2mu}
  <
  7 + \frac{4}{9}
  <
  \sqrt[3]{413 \mspace{1mu}}
\end{equation*}
